\documentclass[11pt]{article}

\usepackage{amsmath,amssymb,amsthm}
\usepackage{mathtools}
\usepackage[round,authoryear]{natbib}
\usepackage{graphicx}
\usepackage{booktabs}
\usepackage{tabularx}
\usepackage{xurl}
\usepackage{hyperref}
\usepackage[margin=1in]{geometry}
\setlength{\emergencystretch}{2em}

\newtheorem{proposition}{Proposition}
\newtheorem{remark}{Remark}
\newtheorem{definition}{Definition}

\newcommand{\bbP}{\mathbb{P}}
\newcommand{\bbQ}{\mathbb{Q}}
\newcommand{\bbE}{\mathbb{E}}
\newcommand{\R}{\mathbb{R}}
\newcommand{\dd}{\mathrm{d}}
\newcommand{\E}{\mathbb{E}}
\newcommand{\var}{\operatorname{Var}}
\newcommand{\tr}{\operatorname{tr}}
\newcommand{\sigcur}{\sigma_{i}}
\newcommand{\signext}{\sigma_{i+1}}

\title{Two Different Objects:\\
The DiffusionPDE Guided-ODE Baseline is Not the Girsanov-Corrected SMC Target}
\author{}
\date{\today}

\begin{document}
\maketitle

\begin{abstract}
This repository contains two sampling procedures that are often compared as if they estimated the
same quantity. The first is the guided probability-flow ODE (PF-ODE) of
\citet{huang2024diffusionpde} --- a deterministic denoising trajectory with a post-hoc
loss-gradient correction --- which is the baseline in \texttt{scripts/generate\_*.py}. The second is
the Girsanov-corrected sequential Monte Carlo (SMC) construction of the companion note
\citep{note1}, implemented on Burgers in \texttt{scripts/generate\_burgers\_gem.py}. This note makes
the relationship between the two precise and shows that they are not different discretizations of a
common target: the SMC framework targets the posterior $p(x_0\mid y)\propto p(y\mid x_0)p(x_0)$,
whereas the baseline is a deterministic map whose implicit target, when read as a guided PF-ODE, is
a schedule-dependent cold reconstruction rather than the posterior, obtained by adding a
constant multiple of the loss gradient to the state after each ODE step. We give (i) the principled
guided PF-ODE that a twist induces, and show the schedule-dependent coefficient the baseline omits;
(ii) an exact side-by-side of the update equations; (iii) a three-way disambiguation of the
baseline, the deterministic limit of the guided SDE, and the full SMC sampler; and (iv) two
implementation-level caveats that matter if the baseline convention is imported into the SMC weight:
flat guidance invalidates the closed-form Gaussian kernel ratio, and the PDE-residual term used as
the twist is not a likelihood. The note reports the modeling issues; it makes no code changes and
takes no position on which method is preferable for the baseline's own reconstruction goal.
\end{abstract}

\section{Introduction and scope}\label{sec:intro}

The upstream code of \citet{huang2024diffusionpde}, retained here, solves forward and inverse PDE
problems under partial observation by denoising a Gaussian field with a probability-flow ODE and, at
each step, subtracting gradients of a sparse-observation loss and a PDE-residual loss
\citep[Algorithm~1]{huang2024diffusionpde}. The paper describes this as guided diffusion and, via a
diffusion-posterior-sampling (DPS) approximation \citep{chung2023dps}, writes the guiding term as a
correction to the score \citep[Eq.~(8)]{huang2024diffusionpde}. The SMC extension in this repository
\citep{note1,note2} instead treats guided diffusion as \emph{importance sampling on path space},
with a Girsanov/twisted correction that makes the sampler consistent for the posterior, and a
Hessian-free potential representation \citep{note2,note3}. The two are routinely placed on the same
axis and scored with the same metric (relative $L^2$ error to the ground-truth field), which
suggests they approximate the same object. They do not.

The purpose of this note is to state precisely what each construction computes and where the
frameworks genuinely diverge. We restrict attention to what is actually written in the released
code; where the paper and the code differ, we say so, since that difference is itself one of the
mechanisms by which the baseline departs from any posterior interpretation. We emphasize at the
outset that the baseline is a \emph{reconstruction} procedure and the paper frames it as such
(``fill in the missing information and solve a PDE by modeling the joint distribution''
\citealp{huang2024diffusionpde}); the observation that it is not a posterior sampler is not a
criticism of its stated goal, only a statement that it cannot be used to validate an SMC
construction of that posterior.

\paragraph{Notation and sources.} We use $\sigma$ for the variance-exploding noise level, $D_\theta$
for the pretrained denoiser, $L$ for a terminal likelihood, and $\ell$ for a log-twist. File
references are to paths in this repository.

\section{The Girsanov-corrected SMC framework}\label{sec:framework}

We recall the construction of \citet{note1}, which is the target the SMC code is built to estimate.

\subsection{Target and reverse processes}

Let $x_0\in\R^d$ be the quantity of interest and $y$ the observation, with likelihood $p(y\mid x_0)$
and a diffusion-model prior $p(x_0)$. The posterior target is
\begin{equation}\label{eq:posterior}
  p(x_0\mid y)\ \propto\ p(y\mid x_0)\,p(x_0).
\end{equation}
With forward noising time $t\in[0,T]$ and reverse (sampling) time $\tau\coloneqq T-t$, define the
reverse state $\xi_\tau\coloneqq x_{T-\tau}$, so that $\xi_0=x_T$ is pure noise and $\xi_T=x_0$ is
data. The unguided reverse process has path measure $\bbP$ and dynamics
\begin{equation}\label{eq:prior-sde}
  \dd\xi_\tau = a(\xi_\tau,\tau)\,\dd\tau + g(\tau)\,\dd W^{\bbP}_\tau,
  \qquad \Sigma(\tau)\coloneqq g(\tau)g(\tau)^\top .
\end{equation}
Fix a positive, smooth \emph{twist} $\phi_\tau(x)>0$ and put
\begin{equation}\label{eq:twist}
  \ell_\tau(x)\coloneqq\log\phi_\tau(x),
  \qquad b_\tau(x)\coloneqq\nabla_x\ell_\tau(x),
  \qquad H_\tau(x)\coloneqq\nabla_x^2\ell_\tau(x).
\end{equation}
The guided proposal, with path measure $\bbQ^\phi$, augments the drift by $\Sigma b$:
\begin{equation}\label{eq:guided-sde}
  \dd\xi_\tau = \bigl[a(\xi_\tau,\tau)+\Sigma(\tau)\,b_\tau(\xi_\tau)\bigr]\dd\tau
  + g(\tau)\,\dd W^{\bbQ}_\tau .
\end{equation}
The path-space target is $p(y\mid\xi_T)\,\bbP(\dd\xi_{0:T})$, so the path weight is
\begin{equation}\label{eq:path-weight}
  W(\xi_{0:T}) = p(y\mid\xi_T)\,\frac{\dd\bbP}{\dd\bbQ^\phi}(\xi_{0:T}).
\end{equation}
If $\phi_\tau$ equals the exact intermediate likelihood
$p(y\mid\xi_\tau)=\int p(y\mid\xi_T)p(\xi_T\mid\xi_\tau)\dd\xi_T$, then the guided drift is
$\bbP$-harmonic, \eqref{eq:guided-sde} is the Doob $h$-transform of \eqref{eq:prior-sde}, and its
terminal marginal is exactly \eqref{eq:posterior}. A general surrogate $\phi$ makes the proposal
suboptimal but leaves the \emph{validity} of \eqref{eq:path-weight} untouched: the surrogate affects
only weight variance, never the target, provided the terminal correction below is retained or the
surrogate is terminally consistent.

\subsection{Change of measure and sequential weights}

Under standard regularity conditions, Girsanov's theorem gives
\begin{equation}\label{eq:girsanov}
\begin{aligned}
  \log\frac{\dd\bbP}{\dd\bbQ^\phi}
  ={} &-\int_0^T b_\tau(\xi_\tau)^\top g(\tau)\,\dd W^{\bbQ}_\tau \\
  &-\frac12\int_0^T b_\tau(\xi_\tau)^\top \Sigma(\tau) b_\tau(\xi_\tau)\,\dd\tau,
\end{aligned}
\end{equation}
and applying It\^o's formula to $\ell_\tau(\xi_\tau)$ to eliminate the stochastic integral yields the
potential/\allowbreak Feynman--Kac representation
\begin{equation}\label{eq:potential}
  \log\frac{\dd\bbP}{\dd\bbQ^\phi}
  = \ell_0(\xi_0)-\ell_T(\xi_T)+\int_0^T V_\tau(\xi_\tau)\,\dd\tau,
\end{equation}
\begin{equation}\label{eq:Vtau}
  V_\tau(x) = \partial_\tau\ell_\tau(x)
  + a(x,\tau)^\top b_\tau(x)
  + \tfrac12 b_\tau(x)^\top\Sigma(\tau)b_\tau(x)
  + \tfrac12\tr\!\bigl(\Sigma(\tau)H_\tau(x)\bigr).
\end{equation}
Discretizing $[0,T]$ by $0=\tau_0<\dots<\tau_K=T$ with step $k$ advancing $\xi_{k-1}\to\xi_k$,
accumulated diffusion $\Sigma_k$ and realized increment $z_k\sim N(0,\Sigma_k)$, the twist increment
telescopes:
\begin{equation}\label{eq:twist-increment}
  \Delta f_k \coloneqq f_{\tau_k}(\xi_k) - f_{\tau_{k-1}}(\xi_{k-1}),
  \qquad
  f_T(\xi_T) = f_0(\xi_0) + \sum_{k=1}^K \Delta f_k .
\end{equation}
This gives the common sequential form
\begin{equation}\label{eq:logW}
  \log W(\xi_{0:T})
  = f_0(\xi_0) + \sum_{k=1}^K G_k + \bigl[\log p(y\mid\xi_T)-f_T(\xi_T)\bigr],
\end{equation}
with the \emph{same} exact continuous-time increment represented two ways,
\begin{equation}\label{eq:Gstar}
  G_k^\star = \Delta f_k + C_k = A_k,
\end{equation}
where the left-endpoint stochastic (Girsanov) increments are
\begin{equation}\label{eq:Ck}
  C_k = -b_{\tau_{k-1}}(\xi_{k-1})^\top z_k
  - \tfrac12 b_{\tau_{k-1}}(\xi_{k-1})^\top\Sigma_k b_{\tau_{k-1}}(\xi_{k-1}),
\end{equation}
and $A_k=\int_{\tau_{k-1}}^{\tau_k} V_\tau(\xi_\tau)\dd\tau$ is the potential increment. The three
standard choices are
\begin{equation}\label{eq:three-increments}
  G^{\mathrm{PBS}}_k = \Delta f_k,
  \qquad
  G^{\mathrm{Girs}}_k = \Delta f_k + C_k,
  \qquad
  G^{\mathrm{pot}}_k = A_k,
\end{equation}
with pseudo-bootstrap (PBS) obtained by dropping the path-measure correction $C_k$, and the
one-parameter ablation
\begin{equation}\label{eq:lambda}
  G^\lambda_k = \Delta f_k + \lambda C_k, \qquad \lambda\in[0,1].
\end{equation}
The initial weight is $\log w_0 = f_0(\xi_0)$; the terminal correction is
$\log w_K\leftarrow\log w_K+\log p(y\mid\xi_T)-f_T(\xi_T)$; and resampling fires when the effective
sample size drops below half the particle count \citep{chopin2020smc}. Under the exact twist,
$V_\tau\equiv0$ by the backward Kolmogorov equation, hence $G^{\mathrm{Girs}}_k=G^{\mathrm{pot}}_k=0$
and the entire path weight collapses to the boundary factor $f_0(\xi_0)$ --- the zero-variance Doob
ideal against which the machinery is validated \citep{note1,note2}.

For an Euler--Maruyama proposal the kernel ratio is \emph{exactly} \eqref{eq:Ck} at every step size,
because both the unguided kernel $p(\xi_k\mid\xi_{k-1})=N(\xi_k;\mu_p,\Sigma_k)$ and the guided
kernel $q(\xi_k\mid\xi_{k-1})=N(\xi_k;\mu_q,\Sigma_k)$ share covariance $\Sigma_k$ and differ only in
mean, with $\mu_q-\mu_p=\Sigma_k b_{\tau_{k-1}}$. Appendix~\ref{app:kernel} derives the two cases
that matter below. We stress the structural requirement that will reappear in
Section~\ref{sec:caveats}: \emph{the twist must enter the proposal drift}, i.e.\ the guided mean
shift must be $\Sigma_k b$, not $b$.

\section{The DiffusionPDE guided-ODE baseline}\label{sec:baseline}

\subsection{As specified in the paper}

The method of \citet{huang2024diffusionpde} trains one denoiser $D_\theta$ on the joint
(coefficient, solution) field and samples by integrating the probability-flow ODE
\begin{equation}\label{eq:paper-ode}
  \dd x = -\dot\sigma(t)\,\sigma(t)\,\nabla_x\log p\bigl(x;\sigma(t)\bigr)\,\dd t,
  \qquad
  \nabla_x\log p\bigl(x;\sigma(t)\bigr)=\frac{D_\theta(x,\sigma(t))-x}{\sigma(t)^2},
\end{equation}
with guidance approximated \`a la DPS: under $y\mid x\sim N\!\bigl(\mathcal M(x),\delta^2 I\bigr)$,
\begin{equation}\label{eq:paper-dps}
  \nabla_x\log p\bigl(y\mid x_i;\sigma(t_i)\bigr)
  \approx -\frac{1}{\delta^2}\nabla_{x_i}\bigl\|y-\mathcal M\bigl(\hat x_N^i\bigr)\bigr\|_2^2,
  \qquad \hat x_N^i = D_\theta(x_i,\sigma(t_i)),
\end{equation}
so that
\begin{equation}\label{eq:paper-guided-score}
  \nabla_{x_i}\log p(x_i\mid y)
  \approx \nabla_{x_i}\log p(x_i)
  - \zeta_{\mathrm{obs}}\nabla_{x_i}\mathcal L_{\mathrm{obs}}
  - \zeta_{\mathrm{pde}}\nabla_{x_i}\mathcal L_{\mathrm{pde}},
  \qquad \zeta=\delta^{-2},
\end{equation}
with mean-squared losses
\begin{equation}\label{eq:paper-losses}
  \mathcal L_{\mathrm{obs}} = \frac1n\bigl\|y_{\mathrm{obs}}-\hat x_N^i\bigr\|_2^2,
  \qquad
  \mathcal L_{\mathrm{pde}} = \frac1m\bigl\|f(\hat x_N^i)\bigr\|_2^2 .
\end{equation}
Equation~\eqref{eq:paper-guided-score} is a \emph{score} (drift) correction.

\subsection{As implemented}

Algorithm~1 and the code apply the guidance differently. The released implementation in
\path{scripts/generate_burgers.py} (and the siblings \path{generate_darcy.py},
\path{generate_poisson.py}, \path{generate_helmholtz.py}, \path{generate_ns_*.py}) is, at
step $i$, with $\sigma_i\coloneqq\sigma(t_i)$:
\begin{align}
  \hat x_N &= D_\theta(x_i,\sigma_i), & d_i &= \frac{x_i-\hat x_N}{\sigma_i},
  \label{eq:code-euler}\\[2pt]
  x_{i+1} &= x_i + (\sigma_{i+1}-\sigma_i)\,d_i,
  \qquad\text{(Heun corrector if }\sigma_{i+1}\neq0\text{),} \label{eq:code-heun}\\[2pt]
  x_{i+1} &\leftarrow x_{i+1}-\zeta_{\mathrm{obs}}\,\nabla_{x_i}\mathcal L_{\mathrm{obs}}
  -\zeta_{\mathrm{pde}}\,\nabla_{x_i}\mathcal L_{\mathrm{pde}} . \label{eq:code-flat}
\end{align}
Three differences from \S\ref{sec:baseline} are immediate and all matter.

\begin{enumerate}
\item \textbf{Placement.} The guidance is applied to the \emph{state} after the ODE step
\eqref{eq:code-flat}, not to the score inside the ODE as in \eqref{eq:paper-guided-score}. When a
drift correction is pulled inside the discretized ODE it acquires the step factor
$(\sigma_{i+1}-\sigma_i)$; the flat update does not. We quantify this in
Section~\ref{sec:target}.
\item \textbf{Loss normalization.} The code uses unsquared norms,
\begin{equation}\label{eq:code-losses}
  \mathcal L_{\mathrm{obs}} = \frac{\bigl\|m\odot(u-u_{\mathrm{GT}})\bigr\|_2}{128\cdot 5},
  \qquad
  \mathcal L_{\mathrm{pde}} = \frac{\bigl\|u_t+uu_x-0.01u_{xx}\bigr\|_2}{128\cdot 128},
\end{equation}
with a sensor mask $m$, not the squared MSE \eqref{eq:paper-losses}. Hence the implied twist is not
Gaussian: $e^{-\zeta\|\cdot\|_2}$ is not a Gaussian likelihood in any noise model.
\item \textbf{Scheduling.} The code uses a two-phase rule (observation term only, at full weight,
for the first $80\%$ of steps, then observation weight divided by ten plus the PDE term), whereas
Algorithm~1 applies both terms at every step. The scheduler is a heuristic anneal of
$\zeta_{\mathrm{obs}},\zeta_{\mathrm{pde}}$ along the trajectory.
\end{enumerate}

The paper's own text acknowledges the anneal: the tuned weights ``range from $0.1$ to $10$'' and
``a significant weight can swiftly decrease the error, but an excessively large weight might cause
oversaturation and artifacts'' \citep{huang2024diffusionpde}. This is a reconstruction-oriented
tuning, not a likelihood-derived schedule.

\section{Is the baseline the probability-flow ODE of a twisted marginal?}\label{sec:target}

We now ask the sharp question: does \eqref{eq:code-euler}--\eqref{eq:code-flat}, for some twist
$\phi$, integrate the PF-ODE that transports the tilted marginal
$\pi_\sigma\propto\phi_\sigma p_\sigma$? If so, then in the exact-ODE limit the terminal law would
be $\pi_T\propto \phi_T p_T$, and the baseline would be a deterministic sampler of that twisted
distribution. We show it is not, for any fixed $\phi$.

\subsection{The principled guided PF-ODE}

For a variance-exploding model $x_\sigma=x_0+\sigma z$, the heat equation
$\partial_\sigma p_\sigma=\sigma\Delta p_\sigma$ is the continuity equation
$\partial_\sigma p_\sigma+\nabla\!\cdot\!(p_\sigma v_\sigma)=0$ with velocity
\begin{equation}\label{eq:pfode}
  v_\sigma(x) = -\sigma\,\nabla_x\log p_\sigma(x).
\end{equation}
By Tweedie's identity $\nabla_x\log p_\sigma=(D_\theta(x,\sigma)-x)/\sigma^2$, so
$v_\sigma=(x-D_\theta)/\sigma=d_i$, confirming that \eqref{eq:code-euler} is a forward-Euler step of
$v_\sigma$ in the $\sigma$ variable. For the tilted marginal $\pi_\sigma\propto\phi_\sigma p_\sigma$
the same continuity argument gives the guided PF-ODE
\begin{equation}\label{eq:guided-pfode}
  v^\pi_\sigma(x) = -\sigma\,\nabla_x\log\pi_\sigma(x)
  = v_\sigma(x) - \sigma\,b_\sigma(x),
  \qquad b_\sigma=\nabla_x\log\phi_\sigma .
\end{equation}
Integrating \eqref{eq:guided-pfode} exactly from $\sigma_{\max}$ down to $0$ transports the noise
marginal to $\pi_0\propto\phi_0\,p_0$; with the terminally-consistent exact twist
$\phi_0=p(y\mid\cdot)$ this is the posterior \eqref{eq:posterior}. This is the deterministic
counterpart of the SMC construction, and it is the idealized object the baseline would have to
approximate. (A jitter+ODE variant that makes this exact transport statement robust and connects it
to the Feynman--Kac weighting is the subject of \citealp{note3}.)

\subsection{What the baseline actually integrates}

Take the DPS surrogate $\phi_\sigma(x)=N\!\bigl(y;\mathcal M(D_\theta(x,\sigma)),\delta^2\bigr)$
(up to normalization), so that
\begin{equation}\label{eq:dps-twist}
  \ell_\sigma(x) = \log\phi_\sigma(x)
  \simeq -\frac{1}{2\delta^2}\bigl\|y-\mathcal M(D_\theta(x,\sigma))\bigr\|_2^2,
  \qquad
  b_\sigma(x) = -\frac{1}{2\delta^2}\nabla_x\bigl\|\cdot\bigr\|_2^2 .
\end{equation}
The principled guided state update from $\sigma_i$ to $\sigma_{i+1}$ is obtained by inserting
\eqref{eq:dps-twist} into \eqref{eq:guided-pfode}:
\begin{equation}\label{eq:principled-update}
  x_{i+1} = x_i + (\sigma_{i+1}-\sigma_i)\bigl(d_i-\sigma_i b_{\sigma_i}(x_i)\bigr)
  = x_i + (\sigma_{i+1}-\sigma_i)d_i
  - \frac{\Delta\sigma_i\,\sigma_i}{2\delta^2}
  \nabla_{x_i}\bigl\|\cdot\bigr\|_2^2,
\end{equation}
with $\Delta\sigma_i\coloneqq\sigma_i-\sigma_{i+1}>0$; the last term is a \emph{descent} direction in
the residual. Thus, if the DPS likelihood is used consistently, its guidance must carry the
schedule-dependent coefficient
\begin{equation}\label{eq:zeta-eff}
  \zeta_{\mathrm{eff}}(\sigma_i) = \frac{\Delta\sigma_i\,\sigma_i}{2\delta^2}.
\end{equation}
The baseline \eqref{eq:code-flat} instead subtracts a \emph{constant} $\zeta$ times the residual
gradient. Two complementary readings make precise how this differs from the guided PF-ODE.

\paragraph{Reading 1: intended twist, wrong coefficient.}
If $\zeta$ is intended as the DPS weight (roughly $1/(2\delta^2)$), then the constant replaces the
schedule-dependent \eqref{eq:zeta-eff}. The ratio
\begin{equation}\label{eq:ratio}
  \frac{\zeta_{\mathrm{eff}}(\sigma_i)}{\zeta}
  = \frac{\Delta\sigma_i\,\sigma_i}{2\delta^2\zeta}
\end{equation}
varies by orders of magnitude along a $\rho$-spaced schedule and tends to $0$ as $\sigma_i\to0$.
Hence no constant $\zeta$ reproduces the DPS guided PF-ODE at more than one noise level.

\paragraph{Reading 2: reverse-engineering the target.}
One may instead ask which twist would make the flat update a \emph{consistent} guided PF-ODE. The PF
guidance term is $\Delta\sigma_i\sigma_i\,b$, so equating it to $-\zeta\nabla\mathcal L$ gives
\begin{equation}\label{eq:base-twist}
  b_{\sigma_i}(x) = -c_i\,\nabla_x\mathcal L(x),
  \qquad
  c_i\coloneqq\frac{\zeta}{\Delta\sigma_i\,\sigma_i},
\end{equation}
which is a gradient, namely that of the schedule-tailored twist
$\phi^{\mathrm{base}}_{\sigma_i}(x)=\exp\!\bigl(-c_i\mathcal L(x)\bigr)$. In other words, the
baseline \emph{is} an Euler step of a guided PF-ODE --- but for $\phi^{\mathrm{base}}$, not for the
DPS likelihood. This yields a sharper statement.

\begin{proposition}[The baseline targets a schedule-dependent, cold reconstruction]\label{prop:no-twist}
Run for $K$ steps, the baseline is the Euler discretization of the guided PF-ODE with twist
$\phi^{\mathrm{base}}_{\sigma_i}(x)=\exp(-c_i\mathcal L(x))$ and terminal inverse temperature
\begin{equation}
  c_K = \frac{\zeta}{\Delta\sigma_{K-1}\,\sigma_{K-1}} .
\end{equation}
In the exact-flow limit (if the ODE with this twist were integrated exactly) its terminal law is
therefore $\pi^{\mathrm{base}}\propto\exp(-c_K\mathcal L(x))\,p_0(x)$, not the posterior
\eqref{eq:posterior}. Since $\Delta\sigma_{K-1}\to0$ as $K\to\infty$ for a fixed schedule endpoint,
the target becomes infinitely cold and concentrates on minimizers of $\mathcal L$. In particular it
depends on the number of steps $K$: it is not discretization-invariant, and no fixed twist reproduces
it across schedules.
\end{proposition}

\begin{remark}
The flat placement is what produces the schedule-tailored twist \eqref{eq:base-twist}. If the same
guidance were instead moved inside the ODE as in \eqref{eq:paper-guided-score}, its coefficient
would be $\zeta_{\mathrm{eff}}(\sigma_i)$ of \eqref{eq:zeta-eff}; the flat update replaces this with
$\zeta$, which is equivalent to the tailored twist above. The paper is thus internally inconsistent
between Eq.~(8) (a score/drift correction) and Algorithm~1 (a flat state update); the code implements
the latter. The SMC framework's kernel ratio (Section~\ref{sec:caveats}) assumes guidance lives in
the drift, i.e.\ an update of the form
$x_{i+1}=x_i+\Sigma_i(s_{\sigma_i}+b_{\sigma_i})+z_i$, which is the step-scaled convention, not the
flat one.
\end{remark}

\begin{remark}[Early vs.\ late weighting]
The mismatch of Reading 1 is not a global rescaling: relative to \eqref{eq:zeta-eff} the baseline
under-weights guidance at large $\sigma$ and over-weights it at small $\sigma$, where the required
coefficient vanishes but the constant does not. This is exactly the signature of the effective
inverse temperature $c_i$ of \eqref{eq:base-twist} growing without bound as the trajectory denoises,
which is why the baseline behaves as a penalized reconstruction rather than a posterior sampler.
\end{remark}


\section{Structural comparison}\label{sec:table}

Table~\ref{tab:compare} summarizes the two constructions. The rows are chosen so that every
difference is a property of the algorithm, not of implementation detail.

\begin{table}[ht]
\centering
\small
\caption{The Girsanov-corrected SMC framework versus the DiffusionPDE guided-ODE baseline.}
\label{tab:compare}
\begin{tabularx}{\textwidth}{@{}l>{\raggedright\arraybackslash}X>{\raggedright\arraybackslash}X@{}}
\toprule
 & \textbf{Girsanov-corrected SMC} \citep{note1} & \textbf{DiffusionPDE baseline} \citep{huang2024diffusionpde} \\
\midrule
Target & Posterior $p(x_0\mid y)\propto p(y\mid x_0)p(x_0)$; exact up to discretization and terminal correction & Implicit target $\propto\exp(-c_K\mathcal L)p_0$ with $c_K=\zeta/(\Delta\sigma_{K-1}\sigma_{K-1})\to\infty$; a schedule-dependent reconstruction, not the posterior \\
Dynamics & Stochastic reverse SDE \eqref{eq:guided-sde} & Deterministic PF-ODE \eqref{eq:paper-ode}, Heun \\
Guidance placement & Drift: $x_{i+1}=x_i+\Sigma_i(s+b)+z_i$ & Flat state update after the step, with or without a constant schedule \\
Guidance scaling & $\Sigma_i b=\delta_i b$ (step-scaled) & $-\zeta\nabla\mathcal L$, $\zeta$ constant \\
Guidance magnitude source & Twist gradient $b=\nabla\log\phi$, with the PF coefficient $\sigma$ & Tuned $\zeta$; loss is an unsquared residual norm \\
Inference object & Weighted particle population; approximates a distribution & A single deterministic field per noise draw \\
Weights / resampling & Girsanov/PBS/potential weights; systematic resampling at ESS $<N/2$ & None \\
Uncertainty & Posterior spread, ESS available & None \\
Consistency & Yes, for the posterior, as $K,N\to\infty$ under regularity & None claimed; not derived from a change of measure \\
Metric & Posterior functionals (mean, credible region, calibration) & Relative $L^2$ to one ground-truth field \\
\bottomrule
\end{tabularx}
\end{table}

\section{Three distinct objects, not two}\label{sec:three}

A frequent source of confusion is that this repository contains a third, intermediate update, the
deterministic limit of the SMC proposal (the \texttt{--no-noise --rho 0} control in
\path{scripts/generate_burgers_gem.py}). With no injected noise and unweighted particles, the
GEM step is
\begin{equation}\label{eq:gem-det}
  x_{i+1} = x_i + \delta_i\bigl(s_{\sigma_i}(x_i)+b_{\sigma_i}(x_i)\bigr),
  \qquad
  \delta_i = \sigma_i^2-\sigma_{i+1}^2,
  \qquad
  s_{\sigma_i} = \frac{D_\theta(x_i,\sigma_i)-x_i}{\sigma_i^2}.
\end{equation}
This is \emph{not} the baseline. With $b\equiv0$, the ratio of the unguided parts of
\eqref{eq:gem-det} and \eqref{eq:code-euler} is
\begin{equation}\label{eq:gem-vs-baseline}
  \frac{\delta_i\,s_{\sigma_i}}{(\sigma_{i+1}-\sigma_i)\,d_i}
  = \frac{\dfrac{\delta_i(D_\theta-x_i)}{\sigma_i^2}}
         {\dfrac{(\sigma_{i+1}-\sigma_i)(x_i-D_\theta)}{\sigma_i}}
  = \frac{\sigma_i+\sigma_{i+1}}{\sigma_i}\in(1,2],
\end{equation}
using $\delta_i=(\sigma_i-\sigma_{i+1})(\sigma_i+\sigma_{i+1})$. So the deterministic SMC limit
takes a systematically larger score step (up to a factor $2$ at small $\sigma$) than the baseline's
PF-Euler step, even before guidance. The correct comparison is therefore three-way:
\begin{enumerate}
  \item[(A)] the baseline: PF-ODE Heun $+$ flat loss-gradient, constant $\zeta$;
  \item[(B)] the guided SDE's deterministic limit: $\delta_i$-scaled score and guidance, no noise,
        no weights (the natural same-target ablation);
  \item[(C)] the full stochastic GEM $+$ Girsanov SMC.
\end{enumerate}
(A) and (C) differ in both target and dynamics; (B) and (C) differ only in the stochasticity and the
reweighting, and are the pair that can be compared as two estimates of the same posterior.
Using (A) as the reference against which (C)'s relative error is judged conflates target mismatch
with sampling quality.

\section{Two implementation-level caveats in the SMC code}\label{sec:caveats}

The preceding sections concern the baseline. The next two observations concern the SMC code itself,
where the baseline convention or a non-likelihood residual has been imported into places the theory
does not cover. Both are reported here, not fixed.

\subsection{Flat guidance invalidates the closed-form kernel ratio}

The identity \eqref{eq:Ck} is the exact Gaussian kernel log-ratio
$\log\frac{p(\xi_k\mid\xi_{k-1})}{q(\xi_k\mid\xi_{k-1})}$ only when the guidance enters the drift.
Appendix~\ref{app:kernel} derives both cases:
\begin{align}
  \text{drift-embedded }(\mu_q-\mu_p=\Sigma_k b):\quad
  C_k &= -b^\top z_k - \tfrac12 b^\top\Sigma_k b
  = -b^\top z_k - \tfrac12\delta_k\|b\|^2,
  \label{eq:Ck-scaled}\\[2pt]
  \text{flat }(\mu_q-\mu_p=b):\quad
  \widetilde C_k &= -\frac{1}{\delta_k}b^\top z_k - \frac{1}{2\delta_k}\|b\|^2,
  \label{eq:Ck-flat}
\end{align}
in the scalar-covariance case $\Sigma_k=\delta_k I$. These agree only if
$\delta_k=1$, which is not the case for any realistic step. Since $\delta_k$ varies across the
schedule, \eqref{eq:Ck-flat} and \eqref{eq:Ck-scaled} cannot be reconciled by any constant
rescaling of the weight either. Consequently the \texttt{--flat-guidance} control of
\texttt{scripts/generate\_burgers\_gem.py}, which switches the proposal to
$x_{i+1}=x_i+\delta_i s_{\sigma_i}+b_{\sigma_i}+z_i$ while still calling \texttt{girsanov\_increment}
(i.e.\ \eqref{eq:Ck-scaled}), attaches a correction that is not the kernel ratio of the proposal
actually used. The self-check \texttt{smc/scripts\_2/check\_gem\_tds\_real\_model.py} validates
\eqref{eq:Ck-scaled} on the real network and therefore applies only to \texttt{scale\_guidance=True}
(the default); it does not certify the flat mode. Any experiment that seeks to compare against the
baseline convention must first derive the flat kernel ratio \eqref{eq:Ck-flat} (or, equivalently,
stop calling it a Girsanov weight).

\subsection{The PDE residual is not a likelihood}

The SMC target \eqref{eq:posterior} presumes a specified observation likelihood $p(y\mid x_0)$. The
Burgers runner instead defines the twist from the unsquared losses \eqref{eq:code-losses},
\begin{equation}\label{eq:code-ell}
  \ell_\sigma(x) = -\zeta_{\mathrm{obs}}\bigl\|m\odot(u-u_{\mathrm{GT}})\bigr\|_2/(128\cdot5)
  -\zeta_{\mathrm{pde}}\bigl\|u_t+uu_x-\nu u_{xx}\bigr\|_2/(128\cdot128),
\end{equation}
so that the effective target is a \emph{tempered surrogate}
$p(x_0)\exp(-\zeta_{\mathrm{obs}}\|\cdot\|_2-\zeta_{\mathrm{pde}}\|\cdot\|_2)$ rather than
$p(x_0)p(y\mid x_0)$. Two consequences follow.

First, the terminal correction of \eqref{eq:logW} is formally required for exactness whenever
$\phi_0\neq p(y\mid\cdot)$. The runner omits it and instead asserts terminal consistency ``by
construction'', on the grounds that the same loss is the only notion of fit available at every step
including the last (\path{scripts/generate_burgers_gem.py}, comments around the boundary-weight
initialization). That is a modeling choice, not a theorem: it fixes the target to the surrogate
\eqref{eq:code-ell}, not to the Bayesian posterior of any stated observation model. If a genuine
likelihood is intended, either $L(x_0)$ must be specified and the terminal correction retained, or
the surrogate must be shown terminally consistent relative to that $L$.

Second, because the residual is included in $\ell$ with a fixed $\zeta_{\mathrm{pde}}$, the ``PDE
constraint'' is a soft prior whose weight is not derived from a noise model. Combined with the
two-phase schedule, the effective target drifts over the trajectory. This is defensible as a
regularized posterior, but it should be stated as such, and it means the SMC relative-error metric
does not measure fidelity to the Bayesian posterior the \texttt{note\_1} framework is built to estimate.

\section{Implications for experiments}\label{sec:implications}

The analysis suggests the following comparison protocol, which separates the three objects of
Section~\ref{sec:three} and the two notions of ``error''.

\begin{enumerate}
  \item \textbf{Choose a likelihood.} Specify $p(y\mid x_0)$ explicitly (e.g.\ a Gaussian sensor
        model with stated $\delta$) and use it as $L$; keep the PDE residual either out of the twist
        or explicitly declared as a prior. Then the target is unambiguous.
  \item \textbf{Retain the terminal correction} (or prove terminal consistency), so the sampler
        targets the stated posterior rather than a surrogate.
  \item \textbf{Do not mix conventions.} If flat guidance is desired, use
        \eqref{eq:Ck-flat}; if the theory's weight is desired, use drift-embedded guidance.
  \item \textbf{Compare like with like.} Contrast (C) against (B), its own deterministic and
        weighted-off limit, to isolate the value of stochasticity and reweighting. Treat (A) as a
        separate reconstruction method, not as the same target.
  \item \textbf{Report target-appropriate metrics.} For (B) and (C), posterior functionals
        (weighted mean, spread, calibration, ESS); relative $L^2$ to a single ground-truth
        realization is a reconstruction metric and can be minimized by point estimators that need
        not be posterior samples.
\end{enumerate}

\section{Summary}\label{sec:summary}

Table~\ref{tab:map} maps the repository artifacts to the methods and targets discussed here.

\begin{table}[ht]
\centering
\small
\caption{Repository artifacts, methods, and targets.}
\label{tab:map}
\begin{tabularx}{\textwidth}{@{}>{\raggedright\arraybackslash}X>{\raggedright\arraybackslash}X>{\raggedright\arraybackslash}X>{\raggedright\arraybackslash}X@{}}
\toprule
Artifact & Method & Update & Target \\
\midrule
\path{scripts/generate_*.py} & DiffusionPDE baseline (A) & PF-ODE Heun $+$ flat $-\zeta\nabla\mathcal L$ & schedule-dependent cold reconstruction\newline (Proposition~\ref{prop:no-twist}); not the posterior \\
\path{generate_burgers_gem.py} \texttt{--no-noise --rho 0} & Guided-SDE deterministic limit (B) & $x+\delta s+\delta b$ & twisted surrogate (same as C, zero-variance) \\
\path{generate_burgers_gem.py} (default) & GEM $+$ Girsanov SMC (C) & $x+\delta(s+b)+z$, weighted & posterior (up to discretization, terminal correction) \\
\path{scripts_2/weightings/doob_vtau.py} & Potential/Feynman--Kac weighting & $A_k=\int V_\tau\dd\tau$ & same as (C), Hessian-trace form \\
\bottomrule
\end{tabularx}
\end{table}

The central conclusion is a scope statement: the Girsanov-corrected SMC construction estimates the
posterior \eqref{eq:posterior}; the DiffusionPDE baseline does not define or estimate that object.
They are two different things that happen to share a denoiser and a sparse-observation loss, and
they should not be placed on the same target axis. Within the SMC construction, the flat-guidance
control and the use of an unnormalized PDE residual as the twist are the two places where the
implementation currently leaves the theory, and both are documented above rather than changed here.

\appendix

\section{Gaussian kernel ratios for drift-embedded and flat guidance}\label{app:kernel}

Consider Gaussian kernels with common covariance $\Sigma_k$,
\begin{equation}
  p(\xi_k\mid\xi_{k-1}) = N(\xi_k;\mu_p,\Sigma_k),
  \qquad
  q(\xi_k\mid\xi_{k-1}) = N(\xi_k;\mu_q,\Sigma_k),
  \qquad
  \xi_k = \mu_q + \Sigma_k^{1/2}\varepsilon_k ,
\end{equation}
so the realized increment is $z_k=\Sigma_k^{1/2}\varepsilon_k$. Writing
$\|\cdot\|_{\Sigma_k^{-1}}^2$ for the $\Sigma_k^{-1}$-norm,
\begin{equation}
  \log\frac{p}{q}
  = -\frac12\Bigl(\|\xi_k-\mu_p\|^2_{\Sigma_k^{-1}}
  - \|\xi_k-\mu_q\|^2_{\Sigma_k^{-1}}\Bigr).
\end{equation}
\paragraph{Drift-embedded guidance.} If $\mu_q-\mu_p=\Sigma_k b$, then
$\xi_k-\mu_p=\Sigma_k b+z_k$ and $\xi_k-\mu_q=z_k$, so
\begin{equation}
  \log\frac{p}{q}
  = -\frac12\Bigl(\|\Sigma_k b+z_k\|^2_{\Sigma_k^{-1}}-\|z_k\|^2_{\Sigma_k^{-1}}\Bigr)
  = -b^\top z_k-\frac12 b^\top\Sigma_k b,
\end{equation}
which is \eqref{eq:Ck} and, for $\Sigma_k=\delta_k I$, \eqref{eq:Ck-scaled}.
\paragraph{Flat guidance.} If instead $\mu_q-\mu_p=b$, then
$\xi_k-\mu_p=b+z_k$ and
\begin{equation}
  \log\frac{p}{q}
  = -\frac12\Bigl(\|b+z_k\|^2_{\Sigma_k^{-1}}-\|z_k\|^2_{\Sigma_k^{-1}}\Bigr)
  = -b^\top\Sigma_k^{-1}z_k-\frac12 b^\top\Sigma_k^{-1}b,
\end{equation}
which is \eqref{eq:Ck-flat} for $\Sigma_k=\delta_k I$. The two expressions coincide only at
$\Sigma_k=I$; they cannot be related by a constant rescaling when $\Sigma_k$ varies along the
schedule.

\section{ODE coefficient comparison}\label{app:ode}

Let $d_i=(x_i-D_\theta(x_i,\sigma_i))/\sigma_i=v_{\sigma_i}(x_i)$ and
$s_{\sigma_i}=(D_\theta(x_i,\sigma_i)-x_i)/\sigma_i^2=-d_i/\sigma_i$. Then
\begin{equation}
  \text{(A) unguided: } (\sigma_{i+1}-\sigma_i)d_i,
  \qquad
  \text{(B) unguided: } \delta_i s_{\sigma_i}
  = -\frac{\delta_i}{\sigma_i}d_i,
\end{equation}
and the ratio is
\begin{equation}
  \frac{\text{(B)}}{\text{(A)}} = \frac{-\delta_i/\sigma_i}{\sigma_{i+1}-\sigma_i}
  = \frac{\delta_i}{\sigma_i(\sigma_i-\sigma_{i+1})}
  = \frac{\sigma_i+\sigma_{i+1}}{\sigma_i},
\end{equation}
using $\delta_i=(\sigma_i-\sigma_{i+1})(\sigma_i+\sigma_{i+1})$. The two deterministic recursions
therefore differ by a factor in $(1,2]$ even with no guidance at all.

\bibliographystyle{plainnat}
\bibliography{references}

\end{document}
